\Subhead{3.7.2}{Тепловое излучение абсолютно чёрного тела}

Чёрное тело поглощает всё падающее излучение и в равновесии испускает универсальный спектр. В большой полости число электромагнитных мод на единицу объёма в интервале $[\omega,\omega+\dd\omega]$ равно $\omega^2\dd\omega/(\pi^2c^3)$, включая две поляризации. Средняя тепловая энергия моды $\hbar\omega\bar n$ даёт
\begin{equation}
 u_\omega(T)=\frac{\hbar\omega^3}{\pi^2c^3}\frac1{\e^{\hbar\omega/(k_BT)}-1}.
\end{equation}
Нулевая энергия не включена в тепловую плотность излучения. При малых частотах получаем Рэлея–Джинса $u_\omega\simeq k_BT\omega^2/(\pi^2c^3)$, при больших — экспоненциальный закон Вина. Квантовое подавление высоких частот устраняет ультрафиолетовую катастрофу.

Спектральная яркость по частоте $B_\nu=2h\nu^3/c^2[\exp(h\nu/k_BT)-1]^{-1}$. Для изотропного поля поток через отверстие $M=cu/4$; интегрирование даёт $M=\sigma_{\rm SB}T^4$. Положение максимума зависит от того, используется спектр по $\lambda$ или по $\nu$: это разные плотности с якобианом преобразования. Тепловой шум в линии и излучение полости имеют общий источник — равновесное заселение электромагнитных мод.
